% ============================================================
% style/figures.tex
%
% Изображения и раскладки без декоративных рамок.
% ============================================================

% ------------------------------------------------------------
% 1. Простая команда для вставки изображения
% ------------------------------------------------------------

\NewDocumentCommand{\rimg}{O{0.90\linewidth} O{} m}{%
    \includegraphics[
        width=#1,
        keepaspectratio,
        #2
    ]{#3}%
}


% ------------------------------------------------------------
% 2. Базовые стили для rfig
% ------------------------------------------------------------
%
% ВАЖНО:
%   внутри \tcbset нельзя оставлять пустые строки между ключами.
%   Поэтому здесь список ключей записан плотно.

\tcbset{%
    rfigbase/.style={%
        enhanced,%
        width=\linewidth,%
        colback=rbg,%
        colframe=rbg,%
        coltext=rtext,%
        boxrule=0pt,%
        frame hidden,%
        sharp corners,%
        left=0pt,%
        right=0pt,%
        top=0pt,%
        bottom=0pt,%
        boxsep=0pt,%
        before skip=0pt,%
        after skip=0pt,%
        fonttitle=\bfseries,%
        coltitle=rtext,%
        colbacktitle=rbg,%
        center title,%
        toptitle=0pt,%
        bottomtitle=8pt%
    },%
    rfigtype/r/.style={%
        sidebyside,%
        lefthand width=0.56\linewidth,%
        righthand width=0.36\linewidth,%
        sidebyside gap=0.06\linewidth,%
        lower separated=false%
    },%
    rfigtype/l/.style={%
        sidebyside,%
        lefthand width=0.36\linewidth,%
        righthand width=0.56\linewidth,%
        sidebyside gap=0.06\linewidth,%
        lower separated=false%
    },%
    rfigtype/c/.style={%
        sidebyside=false,%
        halign=center,%
        halign lower=center,%
        lower separated=false%
    },%
    rfigtype/one/.style={%
    },%
    rfigtype/grid/.style={%
    }%
}


% ------------------------------------------------------------
% 3. Универсальная среда rfig
% ------------------------------------------------------------
%
% Использование:
%
%   \begin{rfig}{r}{Название}
%       Текст слева.
%       \tcblower
%       Картинка справа.
%   \end{rfig}

\NewTColorBox{rfig}{O{} m m}{%
    rfigbase,%
    rfigtype/#2,%
    title={#3},%
    #1%
}